\centerline{\bf \Huge Окр + $H$. КР}
\centerline{\bf \Huge }


\problem{\textbf{1}} В остроугольном треугольнике $ABC$ через центр $O$ описанной окружности и вершины $B$ и $C$
проведена окружность $Q$. Пусть $OF$ — диаметр окружности $Q$, а $D$ и $E$ — точки ее пересечения
с прямыми $AC$ и $AB$ соответственно. Докажите, что $ADFE$ — параллелограмм.

\problem{\textbf{2}} Окружности $S_1$ и $S_2$ пересекаются в точках $A$ и $B$, центр $O$ окружности $S_1$ лежит на окружности $S_2$. Хорда $AC$ окружности $S_1$ пересекает окружность $S_2$ в точке $D$. Оказалось, что $D$
лежит внутри треугольника $BOC$. Докажите, что отрезки $OD$ и $BC$ перпендикулярны.

\problem{\textbf{3}} Пусть $M$ — основание перпендикуляра, опущенного из вершины $D$ параллелограмма $ABCD$
на диагональ $AC$. Докажите, что перпендикуляры к прямым $AB$ и $BC$, восстановленные в
точках $A$ и $C$ соответственно, пересекутся на прямой $DM$.

\problem{\textbf{4}} Около остроугольного треугольника $ABC$ с различными сторонами описали окружность с
диаметром $BN$. Высота $BH$ пересекает эту окружность в точке $K$.

а) Докажите, что $AN$ = $CK$.

б) Найдите $KN$, если $\angle BAC = 35^{\circ}◦$
, $\angle ACB = 65^{\circ}◦$
, а радиус окружности равен 12.